\documentclass[11pt]{article}
\usepackage[margin=1in]{geometry}
\usepackage{amsmath,amssymb,amsthm,mathtools,booktabs}
\usepackage[T1]{fontenc}
\usepackage{hyperref}
\emergencystretch=2em
\newtheorem{theorem}{Theorem}[section]
\newtheorem{lemma}[theorem]{Lemma}
\newtheorem{proposition}[theorem]{Proposition}
\newtheorem{remark}[theorem]{Remark}
\newcommand{\R}{\mathbb R}
\newcommand{\E}{\mathbf E}
\newcommand{\tvec}{\mathbf t}
\newcommand{\zvec}{\mathbf z}
\newcommand{\evec}{\mathbf e}
\newcommand{\Q}{\mathcal Q_\alpha}
\newcommand{\Cenv}{C_{\mathrm{env}}}
\newcommand{\Cgwc}{C_{\mathrm{gwc}}}
\DeclareMathOperator{\cl}{cl}
\DeclareMathOperator{\Rect}{Rect}
\title{Signed Surface-Envelope Transductively Standardized Conformal Prediction\\
\large Construction, finite-sample proofs, and comparison with the separated shortcut}
\author{Method development note}
\date{September 8, 2026}
\begin{document}
\maketitle
\begin{abstract}
We derive a surface-envelope refinement of globally worst-case (GWC)
transductive standardization. The construction allows signed residuals, retains
the location and scale in a single ratio, and produces coordinatewise residual
thresholds without enumerating a multidimensional grid. Its coverage follows
from deterministic containment of the full conformal set. For nonnegative
scores and the same residual representation, its rectangle is contained in the
old separated mean-row shortcut, with the qualifications and endpoint
conventions stated below. This is weak, samplewise improvement of region size,
not strict improvement on every sample, a comparison of coverage probabilities,
or a uniform comparison with CQHR. We give an unconditional backward-search
algorithm and a proved range on which binary search is valid.
\end{abstract}
\tableofcontents

\section{Setup and the conformal rank}
Condition throughout on an independently trained predictor and any training-only
choices. Let $(X^i,Y^i)_{i=1}^{n+1}$ be exchangeable and let
\[
 \E^i=\big(E_1(X^i,Y^i_1),\ldots,E_d(X^i,Y^i_d)\big)\in\R^d
\]
be coordinatewise nonconformity scores. The first $n$ residual vectors are
observed. Fix a test input $x$ and let $\evec=\E(x,y)$ denote a candidate
residual. Suppose its attainable coordinate values lie in
\[
 D_j(x)=[a_j(x),\infty)\cap\R,\qquad
 D(x)=\prod_{j=1}^d D_j(x),\qquad a_j(x)\in\R\cup\{-\infty\}.
\]
Here $a_j(x)$ is a known bound for the test residual, not an estimated bound on
the distribution. Calibration residuals at other covariates need not exceed
$a_j(x)$. A larger enclosing domain, including $D_j=\R$, is also valid.
Coordinatewise score sublevel sets must be intervals for the final prediction
set in outcome space to be a rectangle; validity itself does not require that.

Write $N=n+1$ and
\begin{equation}\label{eq:quantile}
 r=\lceil N(1-\alpha)\rceil,\qquad
 \Q(v_1,\ldots,v_n)=
 \begin{cases}v_{(r)},&r\le n,\\+\infty,&r=N.\end{cases}
\end{equation}
Use weak acceptance inequalities throughout. No jitter is needed, and ties are
allowed. If $v_i\le w_i$ for every $i$, then $\Q(v)\le\Q(w)$: for each real
$c$, at least as many $v_i$ as $w_i$ are at most $c$, so the ordered values have
the same ordering. This elementary monotonicity is used repeatedly below.

\begin{lemma}[Exchangeable rank]\label{lem:rank}
If $S^1,\ldots,S^N$ are exchangeable real scores, then
$\Pr\{S^N\le\Q(S^1,\ldots,S^n)\}\ge1-\alpha$.
\end{lemma}
\begin{proof}
Randomly break ties using independent continuous auxiliary variables. The rank
of the last score among the $N$ scores is uniform on $\{1,\ldots,N\}$ by
exchangeability. A rank at most $r$ implies the weak acceptance event in the
statement, including when tied values occur. Its probability is therefore at
least $r/N\ge1-\alpha$. If $r=N$, acceptance is automatic by
\eqref{eq:quantile}.
\end{proof}
This is the standard conformal rank argument; see \cite{shafer2008}.
The derivations of the envelope and its comparisons below are supplied here.

\section{Augmented standardization and its self-score link}
For each coordinate define
\begin{equation}\label{eq:observed}
 m_j=\frac1n\sum_{i=1}^n E^i_j,\qquad
 s_j^2=\frac1n\sum_{i=1}^n(E^i_j-m_j)^2.
\end{equation}
The divisor in $s_j^2$ is $n$, whereas the variance of the augmented $N$-point
sample below has divisor $N-1=n$. This distinction fixes all constants in the
link. Initially assume $s_j>0$ for every $j$; Section~\ref{sec:edges} supplies a
valid extension to degenerate calibration samples. Finite population moments
are unnecessary for the finite-sample argument.

For a candidate coordinate $z$, put $u=z-m_j$. The augmented mean is
\begin{equation}\label{eq:mu}
 \mu_j(z)=\frac{nm_j+z}{N}=m_j+\frac{u}{N}.
\end{equation}
Using $\sum_i(E^i_j-m_j)=0$ gives the full variance expansion
\begin{align}
 \sum_{i=1}^n(E^i_j-\mu_j(z))^2+(z-\mu_j(z))^2
 &=\sum_{i=1}^n\left(E^i_j-m_j-\frac{u}{N}\right)^2
      +\left(\frac{nu}{N}\right)^2\notag\\
 &=ns_j^2+\frac{nu^2}{N^2}+\frac{n^2u^2}{N^2}
  =ns_j^2+\frac{nu^2}{N}.\notag
\end{align}
Consequently
\begin{equation}\label{eq:sigma}
 \sigma_j(z)=\sqrt{s_j^2+\frac{(z-m_j)^2}{N}}.
\end{equation}
The basic coordinate transformation, its scalarization, and its candidate
self-score are respectively
\begin{align}
 F_j(t,z)&=\frac{t-\mu_j(z)}{\sigma_j(z)}
 =\frac{t-m_j-(z-m_j)/N}{\sqrt{s_j^2+(z-m_j)^2/N}},\label{eq:F}\\
 \Phi_{\zvec}(\tvec)&=\max_{j\le d}F_j(t_j,z_j),\label{eq:phi}\\
 g_j(e)&=F_j(e,e)
 =\frac{n(e-m_j)/N}{\sqrt{s_j^2+(e-m_j)^2/N}}.\label{eq:g}
\end{align}
Unlike a calibration-only fitted standardization, at the true test residual
these parameters are symmetric functions of all $N$ residual vectors.

\begin{lemma}[Signed inverse]\label{lem:link}
Let $T=n/\sqrt N$. The map $g_j:\R\to(-T,T)$ is strictly increasing, and
for finite $e$,
\begin{equation}\label{eq:linkequiv}
 g_j(e)\le c\quad\Longleftrightarrow\quad e\le\omega_j(c),
\end{equation}
where extended-real endpoints encode empty and unrestricted sublevels, and
\begin{equation}\label{eq:omega}
 \omega_j(c)=
 \begin{cases}
 -\infty,&c\le-T,\\[2pt]
 m_j+\displaystyle\frac{Ns_jc}{\sqrt{n^2-Nc^2}},&-T<c<T,\\[6pt]
 +\infty,&c\ge T.
 \end{cases}
\end{equation}
\end{lemma}
\begin{proof}
Direct differentiation gives
$g'_j(e)=(n/N)s_j^2\{s_j^2+(e-m_j)^2/N\}^{-3/2}>0$.
Its limits at $\pm\infty$ are $\pm T$, never attained at finite $e$.
For $|c|<T$, the equation $g_j(e)=c$, with $u=e-m_j$, yields
\[
 c\sqrt{s_j^2+u^2/N}=nu/N,
 \quad (n^2-Nc^2)u^2=N^2s_j^2c^2,
 \quad \operatorname{sign}(u)=\operatorname{sign}(c).
\]
Thus $u=Ns_jc/\sqrt{n^2-Nc^2}$. Keeping the sign condition is essential:
squaring alone would introduce a spurious root for negative $c$.
\end{proof}
The signed residual sublevel is $D_j(x)\cap(-\infty,\omega_j(c)]$.
There is no $\max\{0,\omega_j(c)\}$ in the signed link. If the upper endpoint
is below a finite $a_j(x)$, the sublevel is empty.

\section{The full conformal reference set}
For a candidate $\evec$ define
\begin{equation}\label{eq:oracle}
 Q_{\evec}=\Q\big(\Phi_{\evec}(\E^1),\ldots,
                          \Phi_{\evec}(\E^n)\big),\qquad
 C_{\mathrm{full}}(x)=\{y:\Phi_{\E(x,y)}(\E(x,y))\le Q_{\E(x,y)}\}.
\end{equation}
Equivalently, $y$ is accepted if $e_j\le\omega_j(Q_{\evec})$ for every $j$.
The threshold depends on the entire candidate, so this reference set need not
be rectangular or cheap to compute.

\begin{theorem}[Full conformal coverage]\label{thm:full}
Under exchangeability, $\Pr\{Y^{n+1}\in C_{\mathrm{full}}(X^{n+1})\}
\ge1-\alpha$.
\end{theorem}
\begin{proof}
At $\evec=\E^{n+1}$, augmenting the calibration sample gives precisely the full
exchangeable sample. The means and sample standard deviations of that sample
are invariant to permutation. Applying the same standardized maximum to each
of the $N$ observations therefore gives exchangeable scalar scores.
Lemma~\ref{lem:rank} proves the assertion. One need not assume that the scores
are independent of the fitted transformation.
\end{proof}

\section{Exact coordinate suprema and signed GWC}
For any nonempty interval $I=[\ell,u]\cap\R$, allowing infinite endpoints,
define
\begin{equation}\label{eq:psi}
 \psi_{j,I}(t)=\sup_{z\in I}F_j(t,z).
\end{equation}
Endpoint limits are included when the endpoint is infinite. Suppressing $j$,
differentiation of \eqref{eq:F} gives
\begin{equation}\label{eq:derivative}
 \frac{\partial F(t,z)}{\partial z}
 =-\frac{s^2+(t-m)(z-m)}{N\{s^2+(z-m)^2/N\}^{3/2}}.
\end{equation}
If $t>m$, the unique critical point is a maximum,
\begin{equation}\label{eq:critical}
 z^*(t)=m-\frac{s^2}{t-m},\qquad
 F(t,z^*(t))=\sqrt{\frac{(t-m)^2}{s^2}+\frac1N}.
\end{equation}
If $t<m$, its critical point is a minimum; if $t=m$, the function is decreasing.
Also $F(t,+\infty)=-1/\sqrt N$ and $F(t,-\infty)=1/\sqrt N$.
Thus the complete closed form is
\begin{equation}\label{eq:psiclosed}
 \psi_{j,I}(t)=\max\left\{
 F_j(t,\ell),F_j(t,u),
 \left.\sqrt{\frac{(t-m_j)^2}{s_j^2}+\frac1N}
 \ \right|\ t>m_j,\ z^*(t)\in I\right\}.
\end{equation}
The conditional entry is omitted when its condition fails, rather than replaced
by zero. Replacing it by zero would incorrectly inflate negative suprema.

The signed GWC score, quantile, and bounds are
\begin{align}
 \Phi_G(\tvec;x)&=\max_j\psi_{j,D_j(x)}(t_j)
                =\sup_{\zvec\in D(x)}\Phi_{\zvec}(\tvec),\label{eq:gwcscore}\\
 Q_G(x)&=\Q\big(\Phi_G(\E^1;x),\ldots,\Phi_G(\E^n;x)\big),\label{eq:gwcq}\\
 U_j(x)&=\omega_j(Q_G(x)),\quad G_j(x)=D_j(x)\cap(-\infty,U_j(x)],
 \quad G(x)=\prod_jG_j(x),\label{eq:G}\\
 \Cgwc(x)&=\{y:\E(x,y)\in G(x)\}.\label{eq:Cgwc}
\end{align}
For a finite lower endpoint $a_j$, \eqref{eq:psiclosed} uses $F_j(t,a_j)$,
$-1/\sqrt N$, and the critical maximum when $t>m_j$ and $z^*(t)\ge a_j$.
For $a_j=-\infty$ it simplifies to
\[
 \psi_{j,\R}(t)=
 \begin{cases}
 \sqrt{(t-m_j)^2/s_j^2+1/N},&t>m_j,\\
 1/\sqrt N,&t\le m_j.
 \end{cases}
\]

\begin{proposition}[GWC containment]\label{prop:gwc}
For every realized calibration sample and test input,
\[C_{\mathrm{full}}(x)\subseteq\Cgwc(x).\]
\end{proposition}
\begin{proof}
For a candidate in $D(x)$, $\Phi_{\evec}(\E^i)\le\Phi_G(\E^i;x)$ for every
$i$, whence $Q_{\evec}\le Q_G$. Full acceptance implies
$g_j(e_j)\le Q_{\evec}\le Q_G$, and Lemma~\ref{lem:link} gives $e_j\le U_j$.
\end{proof}
The dependence of $D(x)$ on the test input causes no difficulty: this is a
pathwise containment argument, not a claim that GWC calibration scores formed
using that input are exchangeable on their own.

\section{From full cells to surface envelopes}
Assume $G(x)$ is nonempty. For each coordinate, partition $G_j$ using its two
endpoints and the distinct calibration residuals strictly between them:
\[
 a_j=\eta_{j,0}<\eta_{j,1}<\cdots<\eta_{j,K_j}=U_j,
 \qquad I_{j,k}=[\eta_{j,k-1},\eta_{j,k}]\cap\R.
\]
There are at most $n+1$ cells. We deliberately use closed cells, a cover with
shared endpoints. This prevents boundary exclusions when residuals have atoms.
A singleton $G_j$ is represented by one singleton cell. Distinct interior
knots avoid zero-width duplicate cells. The proofs allow any finite ordered
interval cover, though Section~\ref{sec:binary} uses the calibration partition.

A full cell $\mathcal I_h=\prod_j I_{j,h_j}$ has exact local score
\[
 \Phi_h^{\mathrm{exact}}(\tvec)=
 \sup_{\zvec\in\mathcal I_h}\Phi_{\zvec}(\tvec)
 =\max_j\psi_{j,I_{j,h_j}}(t_j).
\]
Enumerating all $h$ costs up to $(n+1)^d$ cells. Instead fix only the $j$th
coordinate cell and let the others range over GWC:
\begin{align}
 H_j(\tvec)&=\max_{\ell\ne j}\psi_{\ell,G_\ell}(t_\ell),\label{eq:H}\\
 \overline\Phi_{j,k}(\tvec)&=
 \max\{\psi_{j,I_{j,k}}(t_j),H_j(\tvec)\}
 =\sup_{\substack{z_j\in I_{j,k}\\z_\ell\in G_\ell,\ \ell\ne j}}
       \Phi_{\zvec}(\tvec),\label{eq:surface}\\
 \overline Q_{j,k}&=\Q\big(\overline\Phi_{j,k}(\E^1),\ldots,
                              \overline\Phi_{j,k}(\E^n)\big).\label{eq:surfaceq}
\end{align}
Set $H_1=-\infty$ for $d=1$. Every transformation in these formulas is given
explicitly by \eqref{eq:psiclosed}. For $h_j=k$ we have the pointwise chain
\begin{equation}\label{eq:chain}
 \Phi_{\evec}\le\Phi_h^{\mathrm{exact}}\le\overline\Phi_{j,k}\le\Phi_G
 \quad\hbox{whenever }\evec\in\mathcal I_h.
\end{equation}
The envelope enlarges the candidate domain before taking the supremum. In
general a quantile of pointwise maxima is not the maximum of the corresponding
quantiles. Thus no equality with exact full-cell LWC is asserted.

Define the accepted coordinate piece and its upper endpoint by
\begin{align}
 A_{j,k}&=I_{j,k}\cap(-\infty,\omega_j(\overline Q_{j,k})],\label{eq:piece}\\
 B_{j,k}&=
 \begin{cases}
 \min\{\eta_{j,k},\omega_j(\overline Q_{j,k})\},&A_{j,k}\ne\varnothing,\\
 -\infty,&A_{j,k}=\varnothing.
 \end{cases}\label{eq:B}
\end{align}
For a finite lower endpoint, nonemptiness is exactly
$\omega_j(\overline Q_{j,k})\ge\eta_{j,k-1}$; a value $-\infty$ never
represents a finite residual. In particular, equality retains the singleton at
the cell's left endpoint. Let
\begin{equation}\label{eq:L}
 L_j=\max_{1\le k\le K_j}B_{j,k},\qquad
 R_{\mathrm{env}}(x)=D(x)\cap\prod_j(-\infty,L_j],\qquad
 \Cenv(x)=\{y:\E(x,y)\in R_{\mathrm{env}}(x)\}.
\end{equation}
If any coordinate has no accepted piece, return the empty set. Otherwise
$R_{\mathrm{env}}$ is the coordinatewise lower enclosure of the pieces; it may
fill gaps between them. It is not being identified with the exact local union.

\begin{theorem}[Signed envelope validity]\label{thm:valid}
Pathwise, $C_{\mathrm{full}}(x)\subseteq\Cenv(x)\subseteq\Cgwc(x)$.
Consequently $\Pr\{Y^{n+1}\in\Cenv(X^{n+1})\}\ge1-\alpha$.
\end{theorem}
\begin{proof}
Take $y\in C_{\mathrm{full}}(x)$. Proposition~\ref{prop:gwc} implies that its
residual $\evec$ belongs to $G(x)$. For each $j$, choose a covering cell
$I_{j,k}$ containing $e_j$. Every calibration observation satisfies
$\Phi_{\evec}(\E^i)\le\overline\Phi_{j,k}(\E^i)$. Quantile monotonicity and
full acceptance therefore yield
\[
 g_j(e_j)\le\max_\ell g_\ell(e_\ell)
 \le Q_{\evec}\le\overline Q_{j,k},
 \quad e_j\le\omega_j(\overline Q_{j,k}),
 \quad e_j\in A_{j,k},
 \quad e_j\le B_{j,k}\le L_j.
\]
This proves the first containment simultaneously for all coordinates without a
Bonferroni correction. Each nonempty $A_{j,k}$ lies in $G_j$, so $L_j\le U_j$,
which proves the second. The coverage statement follows from
Theorem~\ref{thm:full}.
\end{proof}

\section{Why it improves the old separated shortcut}\label{sec:dominance}
Here the comparison uses the \emph{same nonnegative calibration residuals},
test score function, level $\alpha$, and calibration partition, with $D_j=[0,\infty)$.
It does not compare signed residuals to capped or shifted residuals. Assume
$s_j>0$. Define the old GWC box $G^{\mathrm{old}}$ and
\begin{equation}\label{eq:oldconst}
 c_\ell=\inf_{z\in G^{\mathrm{old}}_\ell}\frac{\mu_\ell(z)}{\sigma_\ell(z)},
 \qquad r_\ell(h)=\inf_{z\in I^{\mathrm{old}}_{\ell,h}}\sigma_\ell(z).
\end{equation}
For nonnegative residuals $m_\ell\ge0$ and
$G^{\mathrm{old}}_\ell=[0,U^{\mathrm{old}}_\ell]$. Differentiating
$\mu(z)/\sigma(z)$ gives
\[
 \frac{d}{dz}\frac{\mu(z)}{\sigma(z)}
 =\frac{s^2-m(z-m)}{N\sigma(z)^3}.
\]
Its interior critical point, when present, is a maximum. Thus
\[
 c_\ell=\min\left\{\frac{\mu_\ell(0)}{\sigma_\ell(0)},
                      \frac{\mu_\ell(U^{\mathrm{old}}_\ell)}
                           {\sigma_\ell(U^{\mathrm{old}}_\ell)}\right\},
\]
with the second expression interpreted as $1/\sqrt N$ at infinity, and
\[
 r_\ell(h)=\sqrt{s_\ell^2+
   \operatorname{dist}(m_\ell,I^{\mathrm{old}}_{\ell,h})^2/N}.
\]
The old separated score on a full cell is
\[
 \Phi_h^{\mathrm{sep}}(\tvec)
 =\max_\ell\{t_\ell/r_\ell(h_\ell)-c_\ell\},\qquad \tvec\ge0.
\]
If the old box contains the mean, choose each off-coordinate cell to contain
$m_\ell$. Its minimum scale is $s_\ell$. The old mean-row score for surface
$(j,k)$ is consequently
\begin{equation}\label{eq:oldrow}
 \Phi^{\mathrm{old}}_{j,k}(\tvec)=
 \max\left\{\frac{t_j}{r_j(k)}-c_j,
             \max_{\ell\ne j}\left(\frac{t_\ell}{s_\ell}-c_\ell\right)\right\}.
\end{equation}
Its quantile and clipped coordinate pieces define $L_j^{\mathrm{old}}$ exactly
as in \eqref{eq:surfaceq}--\eqref{eq:L}, using the old cells. If the old mean
cell is unavailable, the manuscript shortcut returns its GWC box.

\begin{lemma}[Pointwise score domination]\label{lem:domination}
For $\tvec\ge0$ and the same GWC box and cells,
$\overline\Phi_{j,k}(\tvec)\le\Phi^{\mathrm{old}}_{j,k}(\tvec)$.
\end{lemma}
\begin{proof}
For $z\in I_{j,k}$, $\sigma_j(z)\ge r_j(k)>0$ and
$\mu_j(z)/\sigma_j(z)\ge c_j$. Nonnegativity gives
\[
 F_j(t_j,z)=\frac{t_j}{\sigma_j(z)}-\frac{\mu_j(z)}{\sigma_j(z)}
 \le\frac{t_j}{r_j(k)}-c_j.
\]
Taking the supremum proves the active-coordinate inequality. For $\ell\ne j$
and $z\in G_\ell$, $\sigma_\ell(z)\ge s_\ell$, so similarly
$F_\ell(t_\ell,z)\le t_\ell/s_\ell-c_\ell$. Take the off-coordinate suprema
and then the maximum. Both steps explicitly use $t_\ell\ge0$; reversing its
sign reverses the comparison of the scale terms.
\end{proof}

\begin{theorem}[Uniform weak improvement over the old shortcut]\label{thm:dom}
Under the preceding conditions, with closed-cell and weak-acceptance conventions
for both constructions,
\[
 \Cenv(x)\subseteq C_{\mathrm{old}}(x),\qquad
 L_j\le L_j^{\mathrm{old}}\quad(j\le d)
\]
whenever both coordinate bounds represent nonempty sets. The set containment
also holds when the envelope is empty or when the old shortcut falls back to
GWC. In particular, each outcome interval and the resulting rectangular volume
are weakly smaller, while finite-sample coverage remains at least $1-\alpha$.
\end{theorem}
\begin{proof}
Lemma~\ref{lem:domination} and quantile monotonicity give
$\overline Q_{j,k}\le Q^{\mathrm{old}}_{j,k}$. The signed link is increasing,
so each accepted envelope piece is contained in its corresponding old piece.
Taking coordinatewise suprema yields $L_j\le L_j^{\mathrm{old}}$. Taking score
sublevel sets preserves containment in outcome space. On an old GWC fallback
sample, Theorem~\ref{thm:valid} gives the result directly.
\end{proof}
The proof also works if the envelope's GWC box is a subset of the old box and
its cells refine the old cells: each supremum is then over a smaller domain.
Thus a harmless outward rounding in the old GWC computation does not reverse
the mathematical comparison.

\paragraph{What the statement does and does not compare.}
The old manuscript uses some strict cell tests and assumes distinct residuals;
the present note specifies a closed version to cover atoms as well. A literal
comparison to code with a different tie policy requires numerical checking or
harmonizing those policies; it is not justified by ignoring endpoints.
The repository's old implementation also adds positive numerical jitter.
The mathematical theorem is about the explicitly defined rules, not arbitrary
floating-point perturbations. The paired experiment audit separately compares
the actual implementations. During that audit, capped CQR scores exposed a
specific bug in the former implementation: \texttt{argmax} selected the first
tied zero below the mean, which could select a $[0,0]$ cell not containing the
strictly positive mean. The off-coordinate scale then exceeded $s_\ell$, so
\eqref{eq:oldrow} no longer described that code. Selecting the last value at
or below the mean (right-sided insertion rank) restores a mean-containing cell.
All six saved failing examples satisfy containment after this correction;
additional atom-heavy regression tests cover the same issue. Thus the theorem
compares the mathematical old shortcut, including this necessary tie fix,
and does not claim dominance over the buggy historical executable.
The theorem is a containment result, not a claim
that the envelope has higher coverage than a larger old set. Strict size
improvement need not occur. Exact full-ratio cell LWC can be tighter than the
surface envelope, and CQHR has a different scalarization; neither is dominated
by this argument.

\section{Computation and a justified search rule}
Compute $m,s$ and the GWC scores in $O(nd)$ operations. Sort the distinct
calibration coordinates in $O(dn\log n)$. Precompute
$V_{i\ell}=\psi_{\ell,G_\ell}(E^i_\ell)$ and all
$H_{ij}=\max_{\ell\ne j}V_{i\ell}$ in $O(nd)$ using left and right cumulative
maxima. Each surface quantile then requires $O(n)$ arithmetic and an order
statistic selection, with no factor of $d$ inside the surface search.

\begin{lemma}[Rightmost surviving cell]\label{lem:backward}
For any ordered interval cover, the largest $B_{j,k}$ is the endpoint of the
rightmost nonempty $A_{j,k}$. Thus scanning from $K_j$ downwards and stopping at
the first nonempty piece computes $L_j$ exactly, without a monotone predicate.
\end{lemma}
\begin{proof}
If cell $k$ survives and $h<k$, then
$B_{j,h}\le\eta_{j,h}\le\eta_{j,k-1}\le B_{j,k}$ whenever $h$ survives.
An empty piece cannot increase the maximum. All cells to the right have already
been checked and are empty.
\end{proof}
The worst-case cost is $O(dn^2+dn\log n)$, and the actual search cost is
$O(n\sum_j M_j)$ when $M_j$ cells are inspected. The implementation records
these counts. This proof does not claim that the values $B_{j,k}$ are monotone:
they drop to the empty sentinel after a failure.

\subsection{A range where binary search is proved}\label{sec:binary}
For $n>1$ put $b_j=m_j-s_j/\sqrt{n-1}$. We prove a stronger search fact than
the usual restriction to cells above the empirical mean, but make no
unconditional prefix claim for the cells below $b_j$.

\begin{lemma}[Rank-count predicate above $b_j$]\label{lem:binary}
Consider a calibration-partition cell $I_{j,k}=[\ell,u]$ with finite
$\ell>b_j$ and $u>\ell$. Set $h_i=H_j(\E^i)$. Then
\begin{equation}\label{eq:rankpredicate}
 A_{j,k}\ne\varnothing\quad\Longleftrightarrow\quad
 \#\{i:E^i_j<\ell,\ h_i<g_j(\ell)\}<r.
\end{equation}
Consequently the survival predicate is a true prefix followed by a false suffix
on these cells, and binary search is valid on this part of the partition.
\end{lemma}
\begin{proof}
First $g_j(\ell)>-1/\sqrt N$ is equivalent to $\ell>b_j$ by
\eqref{eq:g}. For fixed $t=\ell$, \eqref{eq:derivative} shows that
$F_j(\ell,z)$ has either no interior maximum or is decreasing for $z\ge\ell$.
Its supremum over $[\ell,\infty)$ is the larger of $g_j(\ell)$ and
$-1/\sqrt N$, hence is $g_j(\ell)$. More explicitly, if $\ell<m_j$ the
interior critical point is a minimum; if $\ell\ge m_j$ the derivative is
negative on this half-line.

If $t<\ell$, then $F_j(t,z)<F_j(\ell,z)$ at every finite $z$, and its limit
at infinity is strictly below $g_j(\ell)$. Continuity, including that limit,
therefore implies $\psi_{j,I_{j,k}}(t)<g_j(\ell)$. If $t=\ell$, the supremum
equals $g_j(\ell)$. If $t>\ell$, evaluating at $z=\ell$ already gives a value
strictly larger than $g_j(\ell)$. It follows that the $i$th surface score is
strictly below $g_j(\ell)$ exactly when both inequalities inside the count in
\eqref{eq:rankpredicate} hold. Its $r$th order statistic is at least
$g_j(\ell)$ exactly when fewer than $r$ scores are strictly below this value.
Link inversion proves the equivalence. As $\ell$ increases, $g_j(\ell)$
increases, so the counted sets are nested and their cardinalities cannot decrease.
\end{proof}
In fact, the proof of the rank predicate uses only $z\ge\ell$ and the left
endpoint; no sign assumption is placed on the calibration residuals.
One can check the top cell first, use binary search on the cells with left
endpoint above $b_j$, and scan backwards below that range only if no upper
cell survives. This has $O(dn\log n)$ search cost when every coordinate boundary
is found in the proved range, with the unconditional backward bound otherwise.
Observed prefix behavior outside the proved range must not be substituted for
a proof. The initial reference implementation uses unconditional backward search
so its correctness does not depend on this optional acceleration.

\subsection{Direct order-statistic localization}\label{sec:ranklocal}
Binary search is not necessary in the certified range. For each coordinate set
\[
 W_{ij}=\max\{E^i_j,\,\omega_j(H_j(\E^i))\},\qquad
 T_j=W_{(r),j},
\]
where $W_{(r),j}$ is the $r$th order statistic, with the same infinite-rank
convention. The inverse is evaluated separately at each off-coordinate score;
it is not the inverse of an aggregate quantile.
\begin{proposition}[One-order-statistic localization]\label{prop:ranklocal}
For a nondegenerate cell with left endpoint $\ell>b_j$,
\[
 A_{j,k}\ne\varnothing\quad\Longleftrightarrow\quad \ell\le T_j.
\]
In particular every cell with left endpoint exceeding
$\max\{b_j,T_j\}$ can be discarded without evaluating a surface quantile.
\end{proposition}
\begin{proof}
The extended inverse convention gives
$h_i<g_j(\ell)\Longleftrightarrow\omega_j(h_i)<\ell$ even when
$h_i$ is outside the open range of $g_j$. Hence
\[
 \{E^i_j<\ell,\ h_i<g_j(\ell)\}
 =\{\max(E^i_j,\omega_j(h_i))<\ell\}=\{W_{ij}<\ell\}.
\]
There are fewer than $r$ such values exactly when $W_{(r),j}\ge\ell$.
Apply Lemma~\ref{lem:binary}. The discarded cells are all in that lemma's
range, so this pruning makes no claim about the remaining lower cells.
\end{proof}
Select $T_j$ in $O(n)$ time and find the last partition left endpoint at or
below $\max\{b_j,T_j\}$ by insertion search in $O(\log n)$ time. Start the
unconditional backward scan there. If this cell has left endpoint above $b_j$,
it is the rightmost survivor and only its surface quantile is needed to compute
the actual clipped endpoint $B_{j,k}$. The endpoint is generally \emph{not}
$T_j$: the order statistic localizes the cell, after which
\eqref{eq:surfaceq}--\eqref{eq:L} still determine its boundary.
If the candidate cell is below the certified range, continue backward normally.
Thus search costs $O(nd+d\log n)$ when all coordinate boundaries lie in the
certified range, after sorting and off-coordinate preprocessing, with the
unconditional worst-case bound otherwise. The optional implementation
\texttt{search="rank"} uses this pruning and an outward numerical guard near
ties. It need not be faster on typical samples: a backward scan that already
checks one cell avoids the extra transformed order statistic.

\section{Outcome intervals, signed improvements, and edge cases}\label{sec:edges}
For absolute residuals $E_j(x,y)=|y-\hat f_j(x)|$, $a_j=0$ and the prediction
interval is $[\hat f_j(x)-L_j,\hat f_j(x)+L_j]$.
For ordered fitted quantiles $q^-_j(x)\le q^+_j(x)$, signed CQR scores satisfy
\begin{align*}
 E_j(x,y)&=\max\{q^-_j(x)-y,\ y-q^+_j(x)\}\\
 &=\left|y-\frac{q^-_j(x)+q^+_j(x)}2\right|
                  -\frac{q^+_j(x)-q^-_j(x)}2,
 \qquad a_j(x)=-\frac{q^+_j(x)-q^-_j(x)}2.
\end{align*}
Consequently
\[
 E_j(x,y)\le L_j
 \Longleftrightarrow
 q^-_j(x)-L_j\le y\le q^+_j(x)+L_j,
 \qquad \operatorname{length}_j=(q^+_j-q^-_j)+2L_j.
\]
A negative $L_j$ shrinks the base interval. A bound below $a_j$ means an empty
interval, not a singleton created by clamping. A bound equal to $a_j$ gives
the singleton at the midpoint. Mixed signs across coordinates are allowed;
GWC inflation in one or all coordinates does not rule out local shrinkage.
The envelope's computations for one coordinate still depend on the other
coordinates through $H_j$, so a general sign-reflection shortcut is not implied.

\paragraph{Small calibration samples.}
If $r=N$, all conformal quantiles are $+\infty$ and the procedure returns the
entire attainable domain. An implementation must preserve this case, not
replace the quantile by the largest observed score.

\paragraph{Zero calibration variance.}
To extend the definition to all finite calibration samples, define the full
conformal reference scale to be the augmented sample standard deviation when
positive, and $1$ when the entire augmented coordinate is constant. This rule
is symmetric, so Theorem~\ref{thm:full} still holds. If any observed $s_j=0$,
the reference implementation returns the whole attainable domain, a valid
superset of that same reference set. Otherwise all proofs above apply.
This conservative fallback avoids division by zero; it can be costly for
heavily capped scores. The positive-score dominance theorem excludes these
samples because the unregularized old shortcut also lacks a positive scale.
A practical regularized variant could instead add a predetermined positive
ridge to every augmented variance symmetrically, but its comparison must then
be made to an old shortcut with matching regularization.

\paragraph{Infinity, empty boxes, and ties.}
If GWC is empty, return empty immediately. Infinite cell endpoints use the
limits in \eqref{eq:psiclosed}. An infinite upper coordinate makes volume
infinite unless some coordinate interval is empty or has zero length, in which
case rectangular Lebesgue volume is zero. Closed cells and weak inequalities
are maintained across calibration ties and candidate ties. The finite-sample
coverage guarantee is marginal over the calibration/test experiment; it is not
a guarantee conditional on the realized calibration sample or on every $x$.

\section{Reproducibility and comparison protocol}
The implementation is \path{utility/envelope.py}. The numerical formula and
containment audit is \path{envelope_method/verify.py}.
Its machine-readable results are in \path{envelope_method/verification.json}.
Experimental workloads,
their configurations, source data hashes, trial data, and completion states
are recorded under \texttt{envelope\_method/}. Every synthetic empirical trial
generates fresh training, calibration, and test observations and fits a new
prediction model. Compared methods share those fitted predictions and residuals
within a trial. Only the real-data studies resplit a fixed observed dataset.
The residual-array experiments used to audit algebra and search rules are
mathematical checks, not estimates of a fitted model's coverage. Full conformal
coverage is established by the proofs, not inferred from Monte Carlo results.

For a nonnegative comparison use the same residuals for both algorithms.
For signed CQR comparisons retain raw residuals and measure outcome-space
lengths and volumes, including empty intervals. Applying the envelope also to
capped and shifted residuals separately identifies the effect of improving the
local construction versus changing the score representation. Base interval
miscoverage $0.1$ means fitted marginal quantiles $0.05$ and $0.95$; it is
different from the simultaneous target miscoverage, even when both are $0.1$.
Use the same base intervals for envelope and CQHR in each such comparison.
No uniform envelope--CQHR ordering is claimed. Store trial-level results and
compute uncertainty over independent trials, not by treating all test points
from the same fitted/calibrated split as independent trials.

\begin{thebibliography}{9}
\bibitem{shafer2008} G. Shafer and V. Vovk (2008).
A Tutorial on Conformal Prediction. \emph{Journal of Machine Learning Research}
9:371--421. \url{https://jmlr.org/papers/v9/shafer08a.html}.
\end{thebibliography}
\end{document}
